% =====================================================================
%  STATEMENT DOCUMENT.  Two pages, self-contained, no bibliography.
%
%  The three branching-program conjectures extracted from main.tex
%  ("Conjectures on Finding a k-Way Collision with Bounded Memory"):
%  the sharp oblivious sequential curve (BS0+), the oblivious parallel
%  curve at every query width (BP3'), and the C-fold sequential
%  statement for adaptive programs (BS4-C).
%
%  This file states them and nothing else.  Motivation, proofs, the
%  short-output barrier, the discussion of which to attack first, and
%  the neighbouring models that do not exhibit the phenomenon are all
%  in main.tex; none of it is repeated here.  References are given
%  inline, by author and venue, so that no bibliography is needed.
%
%  Built on the shared conjura-conjecture.cls house style -- see
%  latex/conjectures/_template/.  The "coded" theorem/conjecture
%  environments below (cthm/conj, printing a code like "BS0+" as
%  their number) are this document's own, layered on top of the
%  class; they don't collide with the class's own conjecture/
%  definition/theorem environments, which this file doesn't use.
% =====================================================================
\documentclass{conjura-conjecture}

% Coded statements print their code as their number.
\theoremstyle{cjplain}
\newtheorem{innerthm}{Theorem}
\newtheorem{innerconj}{Conjecture}
\newenvironment{cthm}[1]%
  {\renewcommand{\theinnerthm}{#1}%
   \def\theHinnerthm{cthm.\arabic{innerthm}}\innerthm}{\endinnerthm}
\newenvironment{conj}[1]%
  {\renewcommand{\theinnerconj}{#1}%
   \def\theHinnerconj{conj.\arabic{innerconj}}\innerconj}{\endinnerconj}

\newcommand{\NN}{[N]}

\runninghead{K-WAY COLLISIONS, BOUNDED MEMORY}

\begin{document}

\cjkicker{CONJURA \textperiodcentered{} OPEN PROBLEM}
\cjtitle{Finding a Three-Way Collision with Bounded Memory}
\cjsubtitle{Three Branching-Program Conjectures}
\cjstatus{Statement: AI-written, not yet formalized.}
\cjcategory{Symmetric-Key Cryptography.}

\vspace{0.4em}

\noindent\textbf{Setting.}
Throughout, $N \geq 2$ and $f \colon \NN \to \NN$ is a uniformly random function, available only
through an oracle answering $x \mapsto f(x)$.  A \emph{$k$-collision} is a set of $k$ distinct
points of $\NN$ with a common image under $f$; a \emph{$3$-collision} is the case $k = 3$, and two
$3$-collisions are \emph{disjoint} if they share no point.  Asymptotic notation
$\tilde O, \tilde\Omega, \tilde\Theta$ suppresses factors polylogarithmic in $N$.  Success
probabilities are over the choice of $f$ and any coins of the algorithm; randomisation need not be
allowed for separately, since the input is already random.

\smallskip
\noindent\textbf{What is known.}
Four facts fix the landscape.

\begin{enumerate}[label=(F\arabic*),ref=F\arabic*,leftmargin=*,itemsep=1pt,topsep=2pt]
\item\label{F1} \emph{The query threshold is $N^{2/3}$.}  Among $q$ queried points the expected
  number of $3$-collisions is $\binom{q}{3}N^{-2}$, so $q = \Theta(N^{2/3})$ queries are necessary
  and sufficient for constant success probability (Suzuki, Tonien, Kurosawa and Toyota, ICISC
  2006).  Quantitatively, an algorithm making $T$ queries and then naming three points succeeds
  with probability $O\bigl((T^{3}+1)N^{-2}\bigr)$.  This holds for \emph{any} algorithm, at any
  memory bound, and no conjecture below contradicts it.
\item\label{F2} \emph{Without a table, $\Theta(N)$ queries suffice.}  Find a $2$-collision
  $\{u,v\}$ by cycle-finding, retain $u$, $v$ and $y = f(u)$, then sample fresh points until one
  maps to $y$.  This uses $O(\log N)$ bits of memory.  Between~\ref{F1} and~\ref{F2} lies a gap of
  $N^{1/3}$, and the conjectures below assert that memory is what closes it.
\item\label{F3} \emph{The interpolating algorithms.}  Joux and Lucks (ASIACRYPT 2009) achieve
  $T = \tilde O(N/S)$ using a table of $S$ collision pairs, for every $S \leq N^{1/3}$; Guo, Nan
  and Yao (ASIACRYPT 2025) find $C$ pairwise disjoint $3$-collisions in $\tilde O(CN/S)$ queries,
  for every $S \leq \tilde O(C^{2/3}N^{1/3})$.  Every conjecture below is tight against one of
  these if true.
\item\label{F4} \emph{At $k=2$ memory is free.}  Pollard's rho method with Brent's cycle detection
  finds a $2$-collision in $\Theta(\sqrt N)$ queries using $O(\log N)$ bits, and that count is
  optimal even with unbounded memory.  Consequently \emph{every} conjecture below is false at
  $k = 2$, and no proof of any of them can be indifferent to the collision order.
\end{enumerate}

\smallskip
\noindent\textbf{Models.}
Both are the branching program of Borodin and Cook (SICOMP 1982), specialised to a random-function
oracle.  The state may encode anything, so one cannot say what a program \emph{knows}.

\smallskip
\noindent\fbox{\begin{minipage}{\dimexpr\textwidth-2\fboxsep-2\fboxrule\relax}
\textbf{Sequential.}  A \emph{branching program} of length $T$ and width $W$ is a DAG on levels
$0,\dots,T$ with at most $W$ nodes per level, in which every node below level $T$ carries a label
$x \in \NN$ and has one outgoing edge for each $y \in \NN$, and every node at level $T$ carries a
label in $\NN^{3}$.  It is run by starting at the source, querying the label of the current node,
following the edge labelled by the answer, and reading the label of the sink reached; it
\emph{succeeds} if that label is a $3$-collision.  Its \emph{space} is $S = \log_{2} W$, counted in
bits.  It is \emph{oblivious} if all nodes at a level carry the same label, so that the query
schedule $w_{1},\dots,w_{T}$ is fixed in advance, and \emph{repetition-free} if no point labels two
levels.  A program asked for $C$ pairwise disjoint $3$-collisions instead of one carries labels in
$\NN^{3C}$ at level $T$; nothing else changes.

\smallskip
\textbf{Parallel.}  A \emph{parallel branching program} of query width $q$ and $r$ rounds is the
same object with levels $0,\dots,r$, each node below level $r$ carrying a label in $\NN^{q}$ and
having one outgoing edge for each answer tuple in $\NN^{q}$.  The total query count is $T = rq$.
Oblivious and repetition-free mean as before, now of $q$-tuples.
\end{minipage}}

\smallskip
\noindent\textbf{The proved base.}
Exactly one lower bound is known in each model, and both are for oblivious programs, where the
multi-pass streaming switching lemma of Dinur (EUROCRYPT 2020) applies to the answer stream.  For
\emph{adaptive} programs nothing memory-sensitive is known at $k=3$: the only available lower bound
is the query bound~\ref{F1}, which mentions no memory at all.

\begin{cthm}{BS0}[oblivious sequential programs]\label{BS0}
Every repetition-free oblivious branching program succeeding with probability at least $\tfrac12$
satisfies $T \cdot \bigl(S + \lceil \log_{2} N \rceil\bigr) = \tilde\Omega(N)$; in particular
$T = \tilde\Omega(N)$ at width $N^{O(1)}$.
\end{cthm}

\begin{cthm}{BP0}[oblivious parallel programs]\label{BP0}
Every repetition-free oblivious parallel branching program of space $S$ and query width $q$
succeeding with probability at least $\tfrac12$ satisfies
$T \cdot \bigl(S + q\lceil\log_{2} N\rceil\bigr) = \tilde\Omega(N)$; in particular
$T = \tilde\Omega(N/q)$ at width $N^{O(1)}$.
\end{cthm}

\smallskip
\noindent\textbf{The conjectures.}
Each sharpens one of these, and each is tight against an algorithm.

\begin{conj}{BS0$^{+}$}[the oblivious sequential curve at $k=3$]\label{BS0plus}
Every repetition-free oblivious branching program succeeding with probability at least $\tfrac12$
satisfies
\[
  T \;=\; \tilde\Omega\bigl(\max\{\, N^{2/3},\; N/\sqrt{S} \,\}\bigr),
\]
\[
  \text{equivalently} \qquad
  T \cdot S^{1/2} \;=\; \tilde\Omega(N) \ \text{ for } S \leq N^{2/3}.
\]
\end{conj}

\noindent
The first clause is~\ref{F1} and carries no content; all of it is in $T \cdot S^{1/2}$.  The bound
is tight: an oblivious program can accumulate $P \leq S$ collision pairs at cost
$\tilde\Theta(PN/S)$ and then hunt a third preimage of one of the $P$ recorded values in
$\tilde\Theta(N/P)$ further queries, both phases non-adaptively; balancing at $P = \sqrt S$ gives
$T = \tilde\Theta(N/\sqrt S)$.  So a factor $\sqrt S$ separates \ref{BS0} from the conjecture
inside its own class.  The two clauses cross at $S = N^{2/3}$, where the curve meets the query
threshold and the problem closes.

\begin{conj}{BP3$'$}[the oblivious parallel curve, every query width]\label{BP3prime}
For every $q$, every repetition-free oblivious parallel branching program of width $N^{O(1)}$
succeeding with probability at least $\tfrac12$ satisfies
\[
  T \;=\; \tilde\Omega\bigl(\min\{\, N,\; N^{2}/q^{2} \,\}\bigr).
\]
\end{conj}

\noindent
This is tight at polynomial width, where two strategies are available with $O(\log N)$ bits between
rounds: one holds a $2$-collision and hunts $q$ fresh points per round, costing $T = \Theta(N)$ at
any $q$ as in~\ref{F2}; the other detects a $3$-collision inside a single batch, which routing
permits once the width exceeds $\binom{q}{3}$, with probability $\approx q^{3}/N^{2}$ per round and
so $T = \tilde\Theta(N^{2}/q^{2})$.  The two cross at $q = \sqrt N$.  Nothing is conjectured beyond
polynomial width, where a larger width both permits a table and permits naming more triples per
batch; already at $q = \sqrt N$ the proved \ref{BP0} gives $T = \tilde\Omega(\sqrt N)$ against a
conjectured truth of $\tilde\Theta(N)$, so a factor $q$ is missing even in the oblivious class.

\begin{conj}{BS4-C}[$C$ disjoint triples, adaptive programs]\label{BS4C}
For every $C \geq 1$, every branching program that outputs $C$ pairwise disjoint $3$-collisions
with probability at least $\tfrac12$ satisfies
\[
  T \;=\; \tilde\Omega\bigl(\max\{\, C^{1/3}N^{2/3},\ CN/S \,\}\bigr).
\]
\end{conj}

\noindent
No obliviousness is assumed, and no bound on the width.  The first clause is the query threshold
for $C$ disjoint triples, by the argument of~\ref{F1} applied $C$ times; the second is the
algorithm of Guo, Nan and Yao in~\ref{F3}, and the two clauses cross exactly where the two regimes
meet, so the conjecture is tight throughout.  At $C = 1$ it is the full sequential statement
$T \cdot S = \tilde\Omega(N)$.

\end{document}
